Multi-agent systems routinely discard their own execution histories: sub-agents
return summaries rather than transcripts, and orchestrators compact their context
to fit a fixed window. What survives is then the only evidence available when the
system must later decide what each agent actually did. We study that conflict
exactly. A record with $K$ labels must simultaneously (i)~permit reconstruction
of the joint effort profile $S \in \cube$ with small expected Hamming error and
(ii)~retain enough about each agent's unilateral contribution that a bounded
payment computed from the record alone rewards effort by at least $\gamma$,
whatever the teammates did.

Our first result removes the payment rule from the problem: a record admits a
bounded $\gamma$-margin payment for agent~$i$ if and only if every unilateral
deviation by~$i$ moves the retained label distribution by at least $\gamma$ in
total variation, minimised over mixtures of teammate profiles. Accountability is
exactly a lower bound on the record's sensitivity to each individual. From it we
obtain the exact one-bit frontier,
$\Dopt(2,\gamma) = \Mn + n\gamma\bigl(\tfrac{n-1}{2} - \Mn\bigr)$, and an exact
identification of the threshold at which accountability becomes free:
$\Kflat(n) = \gtwo(\Qn)$, the 2-domination number of the $n$-cube, attained by an
explicit record with margin $1 - 1/\dmin \ge 1/2$. Hence $\Kflat(7) = 32$ and
$97 \le \Kflat(9) \le 128$, with machine-checkable certificates. Escaping the
tension already costs about $85\%$ of per-agent resolution at $n = 20$.
Evaluating concrete compaction mechanisms---counts, thresholds,
hashes, partial per-agent retention---we find records that are
reconstruction-optimal yet admit no payment rule with positive margin at all.
